\documentclass[10pt,a4paper]{amsart}
\usepackage[margin=19mm]{geometry}
\usepackage{amsmath,amssymb,mathtools}
\usepackage[scr=rsfs,cal=euler]{mathalfa}
\usepackage{xfrac}
\usepackage[unicode,hidelinks,bookmarks=false]{hyperref}
\newcommand{\ie}{i.e.\ }
\newcommand{\cf}{cf.\ }
\newcommand{\resp}{respectively\ }
\newcommand{\Subsection}{\subsection}
\providecommand{\nobreakdash}{\nobreak\hbox{-}}
\setlength{\emergencystretch}{2em}
\multlinegap0pt
\usepackage{bm}
\usepackage{bbm}
\usepackage{dsfont}
\usepackage{xspace}
\usepackage{cancel}
\usepackage{mathtools}
\usepackage{amsthm}
\makeatletter
\@ifundefined{definition}{\newtheorem{definition}{Definition}}{}
\@ifundefined{lemma}{\newtheorem{lemma}[definition]{Lemma}}{}
\@ifundefined{theorem}{\newtheorem{theorem}[definition]{Theorem}}{}
\makeatother
\newcommand{\Acal}{\mathcal{A}} %REDUNDANT
\newcommand{\Bcal}{\mathcal{B}}
\newcommand{\Bm}{{\boldsymbol{m}}} % because \bm already used
\newcommand{\Z}{\mathbb{Z}}
\newcommand{\be}{{\boldsymbol{e}}}
\newcommand{\bk}{{\boldsymbol{k}}}
\newcommand{\bn}{{\boldsymbol{n}}}
\title[Three characterizations of a self-similar aperiodic 2-dimens]{Three characterizations of a self-similar aperiodic 2-dimensional subshift}
\author{S\'ebastien Labb\'e}
\date{Source: arXiv:2012.03892v4 (CC BY 4.0)}
\begin{document}
\maketitle

\noindent\emph{Source and licence.} Three characterizations of a self-similar aperiodic 2-dimensional subshift. Version arXiv:2012.03892v4 (CC BY 4.0), distributed under the Creative Commons Attribution 4.0 International licence, as recorded in the arXiv licence metadata for this exact version and shown on the abstract page https://arxiv.org/abs/2012.03892v4. Full rights notice: licenses/arxiv-2012.03892-CC-BY-4.0.txt.\par\smallskip

\noindent\textit{Editorial scope.} Counted unit: the Theorem whose source label is \texttt{thm:if-markers-desubstitute} in the article (source0.tex lines 2173--2203, source label \texttt{thm:if-markers-desubstitute}), delivered together with the article's own complete original proof of it (source0.tex lines 2205--2321, one \texttt{proof} environment of 117 lines). No mathematics is written, repaired or re-derived: statement and proof are the article's own text line for line. Required context retained uncounted: lem:recognizable (lines 2073--2090). Every definition, display and label of those lines is retained unchanged. Omitted from this package and not redistributed: 1--2072, 2091--2172, 2204--2204, 2322--5529. Nothing inside a retained excerpt is omitted or altered: every retained body is delivered exactly as the article's lines are written, so no omission receipt of any kind accompanies this package. Citations: the retained excerpts carry no citation command at all, so no bibliography entry of the article is transcribed and no reference is redistributed. Reference adaptation: none was needed and none was made; every reference inside the delivered unit resolves inside it, so no unresolved reference or citation is produced. Printed slips kept literal: no printed word, symbol, hypothesis or number of the counted statement or of its proof was altered. Layout and class: the article's own class is \texttt{amsart} with packages including amsmath, amsfonts, latexsym, graphicx, amssymb, amsthm, geometry, url, color, enumerate, enumitem, caption, cite, fontenc; the wrapper is the lane's pinned \texttt{amsart} editorial wrapper at 10pt on A4, which restates the theorem-like environments the retained text needs and adds the article's own macro definitions that the retained text needs (\texttt{\textbackslash Acal}, \texttt{\textbackslash Bcal}, \texttt{\textbackslash Bm}, \texttt{\textbackslash Z}, \texttt{\textbackslash be}, \texttt{\textbackslash bk}, \texttt{\textbackslash bn}), with extra wrapper packages (bm, bbm, dsfont, xspace, cancel, mathtools). The delivered numbering is the unit's own. Assets: no retained span contains a figure environment, a graphics-inclusion command, a PDF-page inclusion, a table or a TikZ picture; no image file is copied into this package, the package declares no asset dependency and no image file is redistributed with it. Licence: version arXiv:2012.03892v4 of this article is distributed under the Creative Commons Attribution 4.0 International licence, as recorded in the arXiv licence metadata for this exact version and shown on the abstract page https://arxiv.org/abs/2012.03892v4. A full rights notice is delivered at licenses/arxiv-2012.03892-CC-BY-4.0.txt. Characters: every retained excerpt is pure ASCII, so the delivered engine, XeLaTeX with its default Latin Modern text font, reports no missing character.
\begin{lemma}\label{lem:recognizable}
    Let $d\geq1$ and $i$ such that $1\leq i\leq d$.
Let $\omega:\Bcal\to\Acal^{\Z^d}$ be a $d$-dimensional morphism such that
the image of letters are letters or dominoes in the direction $\be_i$.
If $\omega|_\Bcal$ is injective and there exists a subset $M\subset \Acal$ such that
\begin{equation}\label{eq:markers-on-right}
    \omega(\Bcal)
    \subseteq (\Acal\setminus M)
    \cup \left((\Acal\setminus M)\odot^i M\right),
\end{equation}
or
\begin{equation}\label{eq:markers-on-left}
    \omega(\Bcal)
    \subseteq (\Acal\setminus M)
    \cup \left(M\odot^i (\Acal\setminus M)\right),
\end{equation}
    then $\omega$ is recognizable in $\Bcal^{\Z^d}$.
\end{lemma}

\begin{theorem}\label{thm:if-markers-desubstitute}
    Let $\Acal$ be an alphabet
    and $X\subset\Acal^{\Z^d}$ be a subshift.
    If there exists a subset
    $M\subset\Acal$ 
    of markers in the direction 
    $\be_i\in\{\be_1,\dots,\be_d\}$,
    then 
\begin{enumerate}%[\rm (i)]
    \item[(i)] (markers on the right) there exists an alphabet $\Bcal_R$,
    a subshift $Y\subset\Bcal_R^{\Z^d}$
    and a $d$-dimensional morphism
    $\omega_R:Y\to X$
    such that 
    \begin{equation*}
        \omega_R(\Bcal_R)\subseteq (\Acal\setminus M)\cup 
        \left((\Acal\setminus M)\odot^i M\right)
    \end{equation*}
    which is recognizable and onto up to a shift and
\item[(ii)] (markers on the left) there exists an alphabet $\Bcal_L$,
    a subshift $Y\subset\Bcal_L^{\Z^d}$
    and a $d$-dimensional morphism
    $\omega_L:Y\to X$
    such that 
    \begin{equation*}
        \omega_L(\Bcal_L)\subseteq (\Acal\setminus M)\cup 
        \left(M\odot^i (\Acal\setminus M)\right)
    \end{equation*}
    which is recognizable and onto up to a shift.
\end{enumerate}
\end{theorem}

\begin{proof}
    We do only the proof of (i) when the markers are on the right, since
    one case can be deduced from the other using symmetry.

    Since $X$ is a subshift, there exists a language $F\subset\Acal^{*^d}$
    such that $X$ is the set of configurations of $\Acal^{\Z^d}$ without any
    occurrence of patterns from $F$. Notice that since $M$ is a set of markers
    in the direction $\be_i$, we may assume $M\odot^i M\subset F$.
    Let $P\subset\Acal\odot^i\Acal$ and $Q\subset\Acal$ be the following sets:
    \begin{align*}
        P &= \left((\Acal\setminus M) \odot^i M\right) \setminus F,\\
        Q &= \left\{u\in\Acal\setminus M\,\middle|\, \text{ there exists }
        v\in\Acal\setminus
            M \text{ such that } u\odot^iv \notin F\right\}.
    \end{align*}
    We choose some ordering of their elements with indices starting from zero:
    \begin{align*}
		P&=\{p_0,\dots, p_{|P|-1}\},\\
		Q&=\{q_0,\dots, q_{|Q|-1}\}.
    \end{align*}
    We construct the alphabet $\Bcal=\left\{0,1,\dots,|Q|+|P|-1\right\}$
    and define the rule $\omega$ by
    \begin{equation}\label{eq:omegaStoT}
    \begin{array}{rccl}
    \omega:&\Bcal & \to & \Acal^{*^d}\\
        &j& \mapsto&
    \begin{cases}
        q_j         & \text{ if } 0\leq j <  |Q|,\\
        p_{j-|Q|} & \text{ if }|Q|\leq j < |Q|+|P|.\\
    \end{cases}
    \end{array}
    \end{equation}

    We want to show that $\omega$ extends to a map from a set of configurations to $X$ which
    is onto up to a shift.
    Let $x\in X$ be a configuration which can be seen as a function
    $x:\Z^d\to\Acal$.
    Consider the set $x^{-1}(M)\subset\Z^d$ of positions of markers
    in~$x$. From the definition of markers in the direction $\be_i$, 
    markers appear in nonadjacent hyperplanes orthogonal to $\be_i$ in the configuration $x$.
    Formally, there exists
    a set $H\subset\Z$ such that $x^{-1}(M)=\Z^{i-1}\times H\times\Z^{d-i}$ and
    $1\notin H-H$.
    Since $1\notin H-H$, there exists a strictly increasing sequence
    $(a_k)_{k\in\Z}$ such that $\Z\setminus H=\{a_k\mid k\in\Z\}$.
    We assume that $a_0=0$ if $0\in\Z\setminus H$ and $a_0=-1$ if $0\in H$
    which makes the sequence $(a_k)_{k\in\Z}$ uniquely defined.

    In order to define the preimage of $x$ under $\omega$, it is convenient to
    represent $x$ fiber by fiber. For every $\Bm\in\Z^{d-1}$, let
    $x_\Bm:\Z\to\Acal$, be the sequence such that 
    \[
        x_{(n_1,\dots,n_{i-1},n_{i+1},\dots,n_d)}(n_i)
        =
        x(n_1,\dots,n_{i-1},n_i,n_{i+1},\dots,n_d)
    \]
    for every $\bn=(n_1,\dots,n_{i-1},n_i,n_{i+1},\dots,n_d)\in\Z^d$.
    For every $\Bm\in\Z^{d-1}$, let $y_\Bm:\Z\to\Bcal$ be defined as
    \begin{equation*}
    \begin{array}{rccl}
        y_\Bm:&\Z & \to & \Bcal\\
        &k & \mapsto &
    \begin{cases}
        j & \text{ if }a_k+1\in H \text{ and } x_\Bm(a_k)=q_j,\\
        j & \text{ if }a_k+1\notin H    \text{ and } x_\Bm(a_k)\odot^i x_\Bm(a_k+1)=p_{j-|Q|}, 
    \end{cases}
    \end{array}
    \end{equation*}
    The function $y_\Bm$ is well-defined.
    Indeed, let $k\in\Acal$ and $\Bm\in\Z^{d-1}$.
    If $a_k+1\in H$, then 
    $x_\Bm(a_k+1)\in\Acal\setminus M$
    and
    $y_\Bm(k)= x_\Bm(a_k)\in Q$.
    Also if $a_k+1\notin H$, then 
    $x_\Bm(a_k)\in\Acal\setminus M$
    and 
    $x_\Bm(a_k+1)\in M$.
    Since $x\in X$, then $x_\Bm(a_k)\odot^i x_\Bm(a_k+1)\notin F$
    and therefore $x_\Bm(a_k)\odot^i x_\Bm(a_k+1)\in P$.
    We define the configuration $y:\Z^d\to\Bcal$ by its fibers constructed above
    \[
        y(n_1,\dots,n_{i-1},n_i,n_{i+1},\dots,n_d)
        =
        y_{(n_1,\dots,n_{i-1},n_{i+1},\dots,n_d)}(n_i)
    \]
    for every $\bn=(n_1,\dots,n_{i-1},n_i,n_{i+1},\dots,n_d)\in\Z^d$.

    We may now finish the proof of surjectivity.
If $0\in H$, 
then the configuration $x$ is exactly the image under $\omega$ of the configuration $y\in\Bcal^{\Z^d}$ that
we constructed: $x=\omega(y)$.
If $0\notin H$, then 
the configuration $x$ is a shift of the image under $\omega$ of the
configuration $y\in\Bcal^{\Z^d}$: $x=\sigma^{\be_i}\omega(y)$.
    Let
    \[
        Y = \omega^{-1}(X) = \{y\in\Bcal^{\Z^d} \mid \omega(y)\in X\}
    \]
    making $\omega:Y\to X$ a continuous map which is onto up to a shift.
    Notice that the subset $Y$ is closed since $\omega$ is continuous and $X$ is closed.
    Moreover, $Y$ is shift-invariant since 
    for all $y\in Y$ and $\bn\in\Z^d$
    there exists $\bk\in\Z^d$ such that $\omega(\sigma^\bn(y))=\sigma^\bk(\omega(y))$
    meaning that $\sigma^\bn(y)\in Y$.
    Thus $Y$ is a subshift.

The function $\omega$ is of the form
\begin{equation*}
    \omega(\Bcal)
    \subseteq (\Acal\setminus M)
    \cup \left((\Acal\setminus M)\odot^i M\right)
\end{equation*}
and its restriction on $\Bcal$ is injective by construction.
    Therefore, we conclude from Lemma~\ref{lem:recognizable} that
    $\omega$ is recognizable in $\Bcal^{\Z^d}$.
\end{proof}

\end{document}
